\documentclass{article}

% packages
\usepackage{amsmath, amsthm, thmtools, amsfonts, amssymb, luacode, catchfile, tikzducks, hyperref, ifthen}
\ifcsname c@kobocompile\endcsname
	\usepackage[a5paper, total={1072pt, 1448pt}, margin=10pt, includeheadfoot]{geometry} % set page margins
\else
	\usepackage[a4paper, margin=50pt, includeheadfoot]{geometry}
\fi
\usepackage[shortlabels]{enumitem}
\usepackage[skip=3pt, indent=0pt]{parskip}

% language
\usepackage[bidi=basic, layout=tabular, provide=*]{babel}
\ifcsname c@english\endcsname
	\babelprovide[main, import]{english}
\else
	\babelprovide[main, import]{hebrew}
	\babelprovide{rl}
\fi
%\babelfont{rm}{Libertinus Serif}
\babelfont{rm}[Renderer=Harfbuzz]{Libertinus Serif}
\babelfont{sf}{Libertinus Sans}
\babelfont{tt}{Libertinus Mono}

% style
\AddToHook{cmd/section/before}{\clearpage}	% Add line break before section
\linespread{1.3}
\setcounter{secnumdepth}{0}		% Remove default number tags from sections, this won't do well with theorems
\AtBeginDocument{\setlength{\belowdisplayskip}{3pt}}
\AtBeginDocument{\setlength{\abovedisplayskip}{3pt}}
\graphicspath{ {../images/} }

% operators
\DeclareMathOperator\cis{cis}
\DeclareMathOperator\Sp{Sp}
\DeclareMathOperator\tr{tr}
\DeclareMathOperator\im{Im}
\DeclareMathOperator\re{Re}
\DeclareMathOperator\diag{diag}
\DeclareMathOperator*\lowlim{\underline{lim}}
\DeclareMathOperator*\uplim{\overline{lim}}
\DeclareMathOperator\rng{rng}
\DeclareMathOperator\Sym{Sym}
\DeclareMathOperator\Arg{Arg}
\DeclareMathOperator\Log{Log}
\DeclareMathOperator\dom{dom}
\DeclareMathOperator\supp{Supp}
\DeclareMathOperator\var{Var}
\DeclareMathOperator\cov{Cov}

% commands
%\renewcommand\qedsymbol{\textbf{מש''ל}}
%\renewcommand\qedsymbol{\fbox{\emoji{lizard}}}
\newcommand{\Aa}[0]{\mathcal{A}}
\newcommand{\Bb}[0]{\mathcal{B}}
\newcommand{\CC}[0]{\mathbb{C}}
\newcommand{\Cc}[0]{\mathcal{C}}
\newcommand{\EE}[0]{\mathbb{E}}
\newcommand{\FF}[0]{\mathbb{F}}
\newcommand{\Ff}[0]{\mathcal{F}}
\newcommand{\Ii}[0]{\mathcal{I}}
\newcommand{\Gg}[0]{\mathcal{G}}
\newcommand{\Ll}[0]{\mathcal{L}}
\newcommand{\Mm}[0]{\mathcal{M}}
\newcommand{\NN}[0]{\mathbb{N}}
\newcommand{\Nn}[0]{\mathcal{N}}
\newcommand{\PP}[0]{\mathbb{P}}
\newcommand{\Pp}[0]{\mathcal{P}}
\newcommand{\QQ}[0]{\mathbb{Q}}
\newcommand{\RR}[0]{\mathbb{R}}
\newcommand{\Rr}[0]{\mathcal{R}}
\newcommand{\Ss}[0]{\mathcal{S}}
\newcommand{\TT}[0]{\mathbb{T}}
\newcommand{\Uu}[0]{\mathcal{U}}
\newcommand{\Vv}[0]{\mathcal{V}}
\newcommand{\Ww}[0]{\mathcal{W}}
\newcommand{\ZZ}[0]{\mathbb{Z}}
\newcommand{\acts}[0]{\circlearrowright}
\newcommand{\explain}[2] {
	\begin{flalign*}
		 && \text{#2} && \text{#1}
	\end{flalign*}
}
\newcommand{\maketitleprint}[0]{ \begin{center}
	%\begin{tikzpicture}[scale=3]
	%	\duck[graduate=gray!20!black, tassel=red!70!black]
	%\end{tikzpicture}	
	\includegraphics[width=6cm]{cover}
\end{center}
}

% theorem commands
\newtheoremstyle{c_remark}
	{}	% Space above
	{}	% Space below
	{}% Body font
	{}	% Indent amount
	{\bfseries}	% Theorem head font
	{}	% Punctuation after theorem head
	{.5em}	% Space after theorem head
	{\thmname{#1}\thmnumber{ #2}\thmnote{ \normalfont{\text{(#3)}}}}	% head content
\newtheoremstyle{c_definition}
	{3pt}	% Space above
	{3pt}	% Space below
	{}% Body font
	{}	% Indent amount
	{\bfseries}	% Theorem head font
	{}	% Punctuation after theorem head
	{.5em}	% Space after theorem head
	{\thmname{#1}\thmnumber{ #2}\thmnote{ \normalfont{\text{(#3)}}}}	% head content
\newtheoremstyle{c_plain}
	{3pt}	% Space above
	{3pt}	% Space below
	{\itshape}% Body font
	{}	% Indent amount
	{\bfseries}	% Theorem head font
	{}	% Punctuation after theorem head
	{.5em}	% Space after theorem head
	{\thmname{#1}\thmnumber{ #2}\thmnote{ \text{(#3)}}}	% head content

\ifcsname c@english\endcsname
	\theoremstyle{plain}
	\newtheorem{theorem}{Theorem}[section]
	\newtheorem{lemma}[theorem]{Lemma}
	\newtheorem{proposition}[theorem]{Proposition}
	\newtheorem*{proposition*}{Proposition}
	%\newtheorem{corollary}[theorem]{אין חלופה עברית}

	\theoremstyle{definition}
	\newtheorem{definition}[theorem]{Definition}
	\newtheorem*{definition*}{Definition}
	\newtheorem{example}{Example}[section]
	\newtheorem{exercise}{Exercise}[section]

	\theoremstyle{remark}
	\newtheorem*{remark}{Remark}
	\newtheorem*{solution}{Solution}
	\newtheorem{conclusion}[theorem]{Conclusion}
	\newtheorem{notation}[theorem]{Notation}
\else
	\theoremstyle{c_plain}
	\newtheorem{theorem}{משפט}[section]
	\newtheorem{lemma}[theorem]{למה}
	\newtheorem{proposition}[theorem]{טענה}
	\newtheorem*{proposition*}{טענה}
	%\newtheorem{corollary}[theorem]{אין חלופה עברית}

	\theoremstyle{c_definition}
	\newtheorem{definition}[theorem]{הגדרה}
	\newtheorem*{definition*}{הגדרה}
	\newtheorem{example}{דוגמה}[section]
	\newtheorem{exercise}{תרגיל}[section]

	\theoremstyle{c_remark}
	\newtheorem*{remark}{הערה}
	\newtheorem*{solution}{פתרון}
	\newtheorem{conclusion}[theorem]{מסקנה}
	\newtheorem{notation}[theorem]{סימון}
\fi

% Questions related commands
\newcounter{question}
\setcounter{question}{1}
\newcounter{sub_question}
\setcounter{sub_question}{1}

\ifcsname c@english\endcsname
	\newcommand{\question}[1][0]{
		\ifthenelse{#1 = 0}{}{\setcounter{question}{#1}}
		\section{Question \arabic{question}}
		\addtocounter{question}{1}
		\setcounter{sub_question}{1}
	}

	\newcommand{\subquestion}[1][0]{
		\ifthenelse{#1 = 0}{}{\setcounter{sub_question}{#1}}
		\subsection{Part \alph{sub_question}}
		\addtocounter{sub_question}{1}
	}
\else
	\newcommand{\question}[1][0]{
		\ifthenelse{#1 = 0}{}{\setcounter{question}{#1}}
		\section{שאלה \arabic{question}}
		\addtocounter{question}{1}
		\setcounter{sub_question}{1}
	}

	\newcommand{\subquestion}[1][0]{
		\ifthenelse{#1 = 0}{}{\setcounter{sub_question}{#1}}
		\subsection{סעיף \localecounter{letters.gershayim}{sub_question}}
		\addtocounter{sub_question}{1}
	}
\fi

% import lua and start of document
\directlua{common = require ('../common')}

\GetEnv{AUTHOR}

% headers
\author{\AUTHOR}
\date\today

\title{פתרון מטלה 08 --- אנליזה פונקציונלית, 80417}

\DeclareMathOperator\sgn{sgn}

\begin{document}
\maketitle
\maketitleprint[teal]

\question[2]
נחשב את מקדמי טור פורייה של הפונקציה האינטגרבילית,
\[
	f(x)
	= \sgn(x) \cdot \frac{\pi - |x|}{2}
\]
עבור,
\[
	\sgn(x)
	= \begin{cases}
		-1 & x < 0 \\
		0 & x = 0 \\
		1 & x > 0
	\end{cases}
\]
\begin{solution}
	נבחין כי הפונקציה $f$ היא אי־זוגית, לכן $a_n = 0$ לכל $n \in \NN \cup \{ 0 \}$.
	מספיק שנחשב את $b_n$ אם כך,
	\begin{align*}
		b_n
		& = \frac{1}{\pi} \int_{-\pi}^{\pi} f(x) \cdot \sin(nx)\ dx \\
		& = \frac{2}{\pi} \int_{0}^{\pi} f(x) \cdot \sin(nx)\ dx \\
		& = \frac{1}{\pi} \int_{0}^{\pi} \sgn(x) (\pi - |x|) \cdot \sin(nx)\ dx \\
		& = \frac{1}{\pi} \int_{0}^{\pi} (\pi - x) \sin(nx)\ dx \\
		& = \frac{1}{\pi} \left( \pi \int_{0}^{\pi} \sin(nx)\ dx - \int_{0}^{\pi} x \sin(nx)\ dx \right) \\
		& = {\left[ -\frac{1}{n} \cos(nx) \right]}_{x = 0}^{x = \pi} - \frac{1}{\pi} {\left[ -x \cdot \frac{1}{n} \cos(nx) \right]}_{x = 0}^{x = \pi} - \frac{1}{n \pi} \int_{0}^{\pi} \cos(nx)\ dx  \\
		& = \mp \frac{1}{n} + \frac{1}{n} \pm \frac{1}{n} - \frac{1}{n^2 \pi} {\left[ \sin(nx) \right]}_{x = 0}^{x = \pi} \\
		& = \frac{1}{n}
	\end{align*}
	ולכן,
	\[
		\sum_{n = 1}^\infty \frac{1}{n} \sin(nx)
	\]
	טור פורייה של $f$.
\end{solution}

\question{}
נגדיר טור ונחשב את סכומו על־ידי שימוש בזהות פרסבל.

\subquestion{}
\[
	\sum_{n = 1}^\infty \frac{1}{n^2}
\]
ונסיק את ערך הטור $\sum_{n = 1}^\infty \frac{1}{{(2n - 1)}^2}$.
\begin{solution}
	נגדיר את הפונקציה $f(x) = x$,
	אז מאי־זוגיות $a_n = 0$ לכל $n$, ונחשב את $b_n$,
	\[
		b_n
		= \frac{1}{\pi} \int_{-\pi}^{\pi} x \sin(nx)\ dx
		= \frac{2}{\pi} \int_{0}^{\pi} x \sin(nx)\ dx
		= \mp \frac{2}{\pi n}
	\]
	ולכן משוויון פרסבל,
	\[
		\frac{\pi}{4} \cdot \frac{1}{\pi} \int_{-\pi}^{\pi} f^2(x)\ dx
		= 0 + \sum_{n = 1}^\infty 0 + \frac{1}{n^2}
		= \sum_{n = 1}^\infty \frac{1}{n^2}
	\]
	אבל,
	\[
		\frac{\pi}{4} \int_{-\pi}^{\pi} f^2(x)\ dx
		= \frac{\pi}{4} \int_{-\pi}^{\pi} x^2\ dx
		= \frac{2 \pi}{12} \pi^3
		= \frac{\pi^2}{6}
	\]

	נבחין כי,
	\[
		\frac{\pi^2}{24}
		= \sum_{n = 1}^\infty \frac{1}{4n^2}
		= \sum_{n = 1}^\infty \frac{1}{{(2n)}^2}
	\]
	ולכן,
	\[
		\sum_{n = 1}^\infty \frac{1}{{(2n - 1)}^2}
		= \sum_{n = 1}^\infty \frac{1}{n^2} - \sum_{n = 1}^\infty \frac{1}{{(2n)}^2}
		= \frac{\pi^2}{6} - \frac{\pi^2}{24}
		= \frac{\pi^2}{8}
	\]
\end{solution}

\subquestion{}
נחשב את הטור,
\[
	\sum_{n = 1}^\infty \frac{1}{{(2n - 1)}^4}
\]
\begin{solution}
	נגדיר $f(x) = |x|$.
	זוהי פונקציה זוגית ולכן $b_n = 0$ לכל $n$.
	נחשב את $a_n$,
	\[
		a_n
		= \frac{1}{\pi} \int_{-\pi}^{\pi} |x| \cos(nx)\ dx
		= \frac{2}{\pi} \int_{0}^{\pi} x \cos(nx)\ dx
		= \frac{2}{\pi n^2} ({(-1)}^n - 1)
	\]
	וכן,
	\[
		a_0
		= \frac{1}{\pi} \int_{-\pi}^{\pi} f(x)\ dx
		= \frac{1}{\pi} \int_{0}^{\pi} 2x\ dx
		= \pi
	\]
	ומשוויון פרסבל,
	\[
		\int_{-\pi}^{\pi} f^2(x)\ dx
		= \frac{\pi^3}{2} + \sum_{n = 1}^\infty \frac{4}{\pi n^4} |{(-1)}^n - 1|
	\]
	כלומר,
	\[
		\sum_{n = 1}^\infty \frac{1}{{(2n - 1)}^4}
		= \sum_{n = 1}^\infty \frac{1}{n^4} |{(-1)}^n - 1|
		= \frac{\pi}{4} \int_{-\pi}^{\pi} f^2(x)\ dx - \frac{\pi^4}{8}
		= \frac{\pi^4}{6} - \frac{\pi^4}{8}
		= \frac{\pi^4}{24}
	\]
\end{solution}

\question{}
נוכיח את מבחן $M$ של ויירשטראס.
נניח ש־$I \subseteq \RR$ קטע, ${\{ f_n \}}_{n = 1}^\infty \subseteq C(I)$ ו־${\{ M_n \}}_{n = 1}^\infty \subseteq \RR$ סדרת חסמים כך שלכל $n \in \NN$ מתקיים $\lVert f_n \rVert_\infty < M_n$.
נראה שהתכנסות הטור $\sum_{n = 1}^\infty M_n$ גוררת שהפונקציה $f$ המוגדרת על־ידי,
\[
	f(x)
	= \sum_{n = 1}^\infty f_n(x)
\]
מוגדרת היטב ב־$I$, רציפה והטור מתכנס אליה במידה שווה.
\begin{proof}
	לכל $N \in \NN$ מתקיים,
	\[
		\sum_{n = 1}^N f_n(x)
		< \sum_{n = 1}^N M_n
	\]
	ישירות מהנתונים ולכל $x \in I$, לכן ממבחן הסנדוויץ' לגבולות ממשיים נקבל ש־$f$ מוגדרת בכל $I$ ושההתכנסות היא במידה שווה (שכן לא הייתה תלות ב־$x$). \\
	נותר להראות ש־$f$ רציפה.
	יהי $\varepsilon > 0$, אז,
	\[
		|f(x) - f(y)|
		\le \left\lvert f(x) - \sum_{n = 1}^N f_n(x) \right\rvert + \left\lvert \sum_{n = 1}^N f_n(x) - \sum_{n = 1}^N f_n(y) \right\rvert + \left\lvert \sum_{n = 1}^N f_n(y) - f(y) \right\rvert
	\]
	מהתכנסות במידה שווה קיים $N \in \NN$ (נבחר מקסימלי) כך שמתקיים,
	\[
		\left\lvert f(x) - \sum_{n = 1}^N f_n(x) \right\rvert
		\le \frac{\varepsilon}{3},
		\qquad
		\left\lvert f(y) - \sum_{n = 1}^N f_n(y) \right\rvert
		\le \frac{\varepsilon}{3}
	\]
	וכן מרציפות $f_n$ לכל $n \le N$ קיים $\delta > 0$ כך שלכל $|x - y| < \delta$ מתקיים,
	\[
		\left\lvert \sum_{n = 1}^N f_n(x) - \sum_{n = 1}^N f_n(y) \right\rvert
		\le \frac{\varepsilon}{3}
	\]
	ולכן $|f(x) - f(y)| \le \varepsilon$ ומצאנו כי $f$ רציפה במידה שווה ובפרט רציפה.
\end{proof}

\question{}
תהי $f \in \tilde{C}[-\pi, \pi]$ כך ש־$f' \in \tilde{C}[-\pi, \pi]$.
נסמן ב־$a_n^f, b_n^f$ את מקדמי טור פורייה של $f$ וב־$a_n', b_n'$ את מקדמי טור פורייה של $f'$.

\subquestion{}
נראה שלכל $n \in \NN$ מתקיים,
\[
	a_n^f
	= - \frac{1}{n} b_n',
	\qquad
	b_n^f
	= - \frac{1}{n} a_n'
\]
\begin{proof}
	TODO
\end{proof}

\end{document}
